% !PRINTED 443
\weberpage{443}
Diese Arbeit wird wesentlich dadurch erleichtert, dass uns aus der Betrachtung der Congruenzgruppe $L_7$ die Grade und Cykeln der Gruppe schon bekannt sind (\S.\,68). Danach giebt es $p^2-1=48$ Elemente $7^{\text{ten}}$ Grades $A_7$, $\frac14 p(p-3)(p+1)=56$ Elemente $3^{\text{ten}}$ Grades $A_3$ und $\frac14 p(p-1)^2=63$ Elemente $4^{\text{ten}}$ oder $2^{\text{ten}}$ Grades $A_4, A_2$ in $G_{168}$.

Die Elemente $A_7$ ordnen sich (mit Zuziehung der identischen Substitution) in acht Cykeln von sieben Gliedern, die $A_3$ in 28 Cykeln von drei Elementen und die $A_4$ und $A_2$ in 21 Cykeln von vier Gliedern. Jeder dieser letzteren Cykeln enth\"alt ein Element $A_2$ und zwei Elemente $A_4$, und wenn wir also zusammenfassen, so erhalten wir:
\[
\begin{aligned}
&48 \text{ Elemente } 7^{\text{ten}} \text{ Grades}\\
&56 \text{ „} \quad 3^{\text{ten}} \text{ „}\\
&42 \text{ „} \quad 4^{\text{ten}} \text{ „}\\
&21 \text{ „} \quad 2^{\text{ten}} \text{ „}.
\end{aligned}
\]
Wenn ein Punkt durch $\nu$ Substitutionen der Gruppe (die identische eingeschlossen) unge\"andert bleibt, so wollen wir einen solchen Punkt einen \gesperrt{$\nu$-z\"ahligen Pol} nennen (\S.\,52).

F\"ur die Feststellung der Pole und Axen gen\"ugt die Betrachtung je eines Cyklus, da aus diesem die anderen durch Transformation abgeleitet sind (\S.\,69).

Nehmen wir die Substitution $7^{\text{ten}}$ Grades $\tau$ [\S.\,115, (4)], so finden wir die Ecken des Coordinatendreiecks als drei getrennte Pole. Alle Potenzen von $\tau$, deren Exponenten nicht durch 7 theilbar sind, haben dieselben drei Pole.

Durch Anwendung der Substitutionen $\chi,\chi^2$ auf die Coordinaten dieser drei Pole gehen dieselben cyklisch in einander \"uber; wendet man aber die Substitutionen $\Theta,\Theta^2,\Theta^3,\omega,\omega\Theta,\omega\Theta^2,\omega\Theta^3$ an, die in \S.\,115, (5), (10), (16), (18) vollst\"andig gebildet sind, so geht das ganze Coordinatendreieck in ein v\"ollig davon verschiedenes \"uber (z.\,B. durch $\omega$ in das Dreieck mit den Eckcoordinaten $\alpha,\beta,\gamma$; $\beta,\gamma,\alpha$; $\gamma,\alpha,\beta$), und daraus ist zu schliessen, dass diese Pole nicht mehr als siebenz\"ahlig sind. Es giebt acht Tripel solcher siebenz\"ahliger Pole, die alle von einander verschieden sind, und jeder von ihnen kann in jeden anderen durch eine Substitution der Gruppe $G_{168}$ transformirt werden.

Wir haben also den Satz:
%CONT>>
